% "Ending the Game": shared by the rulebook and the platform manuals (docs/manual/*).
% Each document sets its own heading above this; \atscreen{} marks where the manual shows a screen.
After placing your \V{spike}, if either of these is true, you may \term{call the \V{nullroute}}:

\begin{itemize}
  \item \term{\V{totalcompromise}}: any \V{hive} has \Rule{endsize} or more \Vp{node}.
  \item \term{\V{lockdown}}: every \V{hive} on the \V{grid} is \V{safe}.
\end{itemize}

Calling it is your choice; you may keep playing instead. If you call it, finish your turn (you may still buy \V{coin}).
Then the carriers \V{nullroute} the region and the game ends.

\atscreen{nullroute}

\begin{terminal}
[!] NULL-ROUTE \Game{region}: \Game{endshout}
\end{terminal}

\begin{note}{\MakeUppercase{\V{stall}}}
If a full round goes by in which no \V{player} can run an \V{exploit} (every \V{kit} is empty or holds only blocked or
\V{patched} \Vp{exploit}), the \V{grid} can never change again. The carriers seal the region and the game ends
at once, without anyone calling it. Go straight to the final \V{cashout}.
\end{note}

\subsection{The final \V{cashout}}
\begin{enumerate}
  \item Pay \V{root} and \V{sudo} bonuses for every \V{hive} still on the \V{grid}, exactly as in a \V{hivejack}.
  \item Every \V{player} sells all their \V{coin} to the \V{bank} at its current price, starting with the smallest \V{hive}.
  \item \VC{flatlined} \V{coin} (from syndicates not on the \V{grid}) is worth nothing.
\end{enumerate}
The \V{player} with the most \Game{currency} wins. Ties share the win.

\atscreen{final}
